%=======================================================================
%  IE 5311-001 (D01) Exploration Problem EP10
%  One page, single spacing, 10 pt, 1-inch margins, as the sheet requires.
%  Content past page 1 is not graded, so the page limit is hard.
%  Sections map onto the parts asked: what the cone is, its dual and
%  characterisation, problems that use it, and conic duality.
%
%  IF IT SPILLS ONTO A SECOND PAGE, delete in this order until it fits.
%    1. the last sentence of section 5
%    2. the n = 50 and n = 500 rows of the table in section 5
%    3. the sentence in section 3 beginning "Logistic regression"
%  Do not cut the derivation in section 2 or the table in section 5:
%  they are the dual cone and the measured result.
%=======================================================================
\documentclass[10pt]{article}

\usepackage[margin=1in]{geometry}
\usepackage{amsmath}
\usepackage{amssymb}
\usepackage{booktabs}
\usepackage{array}

\newcommand{\T}{^{\!\top}}

\pagestyle{empty}
\setlength{\parskip}{0.35em}
\setlength{\parindent}{0pt}

\begin{document}

{\large\bfseries EP10.\quad Exponential cones and their duals}\\[1pt]
{\small Lidivine Kengne \quad\textbar\quad IE 5311-001 (D01), Homework 6}
\vspace{0.3em}

\textbf{1. What it is.}
The exponential cone is the closure of the region under a perspective
exponential,
\[
  K_{\exp} = \operatorname{cl}
  \bigl\{ (x,y,z) \in \mathbb{R}^3 \;:\; y > 0, \;\;
    y\,e^{x/y} \le z \bigr\}
  \;=\;
  \bigl\{ y > 0,\; y e^{x/y} \le z \bigr\}
  \cup
  \bigl\{ (x,0,z) : x \le 0,\; z \ge 0 \bigr\}.
\]
The function $(x,y) \mapsto y e^{x/y}$ is the perspective of $e^{x}$,
hence convex and positively homogeneous, so $K_{\exp}$ is a closed
convex cone. Unlike the power cone it is not self-dual, and unlike the
second-order cone it is not symmetric --- which is why interior-point
solvers need dedicated barriers for it.

\textbf{2. The dual cone, derived.}
Require $\langle (x,y,z),(u,v,w)\rangle \ge 0$ for every
$(x,y,z) \in K_{\exp}$. On the open branch, $z$ is free above
$y e^{x/y}$, so $w \ge 0$ is forced and the binding case is
$z = y e^{x/y}$. Dividing by $y > 0$ and writing $r = x/y$, the
requirement becomes
\[
  u r + v + w e^{r} \;\ge\; 0 \qquad\text{for all } r \in \mathbb{R}.
\]
The left side is convex in $r$; its minimum is at $e^{r} = -u/w$, which
needs $u < 0$, and substituting gives $u\ln(-u/w) + v - u \ge 0$.
Dividing by $u < 0$ reverses the inequality and yields
$\ln(-u/w) \le 1 - v/u$, that is
\[
  K_{\exp}^{*} = \operatorname{cl}\bigl\{ (u,v,w) :
    u < 0, \; w \ge 0, \; -u\,e^{v/u - 1} \le w \bigr\}
  \cup \bigl\{ (0,v,w) : v \ge 0,\; w \ge 0 \bigr\}.
\]
Sampling (\texttt{ep10\_exponential.py}) confirms both directions: of
the pairs drawn from $K_{\exp} \times K^*_{\exp}$ none paired
negatively, and all $300$ sampled points outside the claimed dual
admitted a cone element they paired negatively with. The witness search
has to span $|x|$ up to $10^{6}$: for $u$ just above zero the pairing
turns negative only once $ux$ dominates $vy$, so a narrow search
wrongly suggests such points belong to the dual.

\textbf{3. Where it arises --- including Problem 1 of this homework.}
Anything built from $e^{x}$, $\ln x$ or $x \ln x$. Problem~1 is the
clean case. Since $-p\ln p \ge u \iff p\,e^{u/p} \le 1$, the maximum
entropy program is
\[
  \min_{p,\,u}\; -\sum_{i \in I} u_i
  \quad\text{s.t.}\quad
  (u_i,\, p_i,\, 1) \in K_{\exp},\;\; i \in I,
  \qquad
  \sum_i p_i = 1, \;\; \sum_i i\,p_i = m_1, \;\; \sum_i i^2 p_i = m_2 .
\]
Logistic regression, geometric programming through $\log\sum\exp$, and
relative entropy objectives are the other standard families.

\textbf{4. Conic duality on it.}
Dualising the three moment rows reproduces exactly the result derived
by hand in Problem~1: the optimality conditions force
$p_i \propto e^{-\nu_1 i - \nu_2 i^2}$, and the dual reduces to two
variables. Numerically, for $n = 6$ the conic program returns
$-1.6965357595$ at $(m_1,m_2) = (4,18)$ and $-1.2665057261$ at
$(4.5,21)$, against $-1.6965357595$ and $-1.2665057255$ from
Problem~1, with the distributions agreeing to $1.8\times10^{-5}$.

\textbf{5. The measured trade, which is not the obvious one.}
Problem~1's reduction looks strictly better: two variables against
$2n$ variables and $n$ cones. It is not, once $n$ grows. Residual below
is the largest moment violation of the returned distribution.

\begin{center}\footnotesize
\setlength{\tabcolsep}{7pt}
\begin{tabular}{@{}rrrrrr@{}}
\toprule
$n$ & cone (ms) & 2-var dual (ms) & cone residual
  & dual residual & rescaled dual\\
\midrule
      6 &   4.4 &   6.1 & $3.6\times10^{-9}$  & $1.9\times10^{-7}$
  & $7.3\times10^{-8}$\\
     50 &   6.2 &   5.1 & $1.7\times10^{-9}$  & $5.2\times10^{-6}$
  & $2.5\times10^{-7}$\\
    500 &  18.3 &   7.3 & $1.2\times10^{-10}$ & $8.5\times10^{-2}$
  & $2.9\times10^{-6}$\\
   5000 & 226.2 &   7.6 & $4.1\times10^{-11}$ & $7.0\times10^{1}$
  & $6.6\times10^{-4}$\\
  50000 & ---   &  79.9 & ---                 & $1.2\times10^{3}$
  & $3.4\times10^{-3}$\\
 200000 & ---   & 636.6 & ---                 & $5.0\times10^{3}$
  & $1.8\times10^{-2}$\\
\bottomrule
\end{tabular}
\end{center}

The conic form holds $10^{-11}$ throughout while the two-variable dual
degrades and, past $n = 5\times10^{4}$, fails outright --- quasi-Newton
never leaves $\nu = 0$. The cause is scaling: the constraint functions
$i$ and $i^2$ differ by a factor $n$, so the dual Hessian conditions
badly and the gradient at the origin is of order $n^2$. Rescaling to
$i/n$ and $(i/n)^2$ recovers four to five orders but does not remove
the trend. The lesson is that the small dual is the right
\emph{mathematics} and the conic form is the right \emph{arithmetic}:
the interior-point method's own scaling buys robustness that the
two-variable reduction has to be given by hand.

\end{document}
